%% =====================================================================
%%  B4-T-25-02 -- what B4 contributes to "The Free Lunch"
%%  -------------------------------------------------------------------
%%  LAYER A: APPEND-ONLY. Written once at the entry's write, then left
%%  alone. If the source has more to say later, add a bullet -- do not
%%  rewrite. Material found ONLY here goes in the `unique` slot with
%%  @unique set to yes, which makes it undeletable (rule R10).
%% =====================================================================
%% @kind: source
%% @id: B4-T-25-02
%% @book: B4
%% @topic: T-25-02
%% @locator: ch. 2 (Reciprocation), pp. 43--86 --- the free-sample section (additive)
%% @status: distilled
%% @unique: no
%% @integrated: yes
%% @updated: 2026-09-19

\dossier{B4}{T-25-02}{\bookshort{B4}}{ch. 2 (Reciprocation), pp. 43--86 --- the free-sample section (additive)}

\begin{distilled}
  \item The free sample is both information and gift: the promoter releases the natural indebting force inherent in a gift while appearing only to inform.
  \item The supermarket sample: many find it difficult to accept the cube, return only the toothpick, and walk away --- they buy the product even when they did not especially like it.
  \item A sample is a purchase already begun: the reciprocity rule converts the trial into obligation before any decision has been consciously made.
  \item Personalised free gifts compound the effect elsewhere in the chapter: a veterans' mailing that added unsolicited individualised address labels nearly doubled its response rate (18 to 35 percent); at the sample table it is the attendant's smile the shopper cannot walk away from.
\end{distilled}
